\section{Constructing the ambient periods and their arithmetic}
\label{sec:period-arithmetic}

We now supply the ambient-period and CM-$j$ inputs retained in TL17
and TL21. The input is the literal arithmetic host $H$: its number
fields $K=\mathbb Q(\sqrt{-d})\subset E$, chosen embeddings, conjugation,
and inverse half-different $A=H.\mathrm{AA0}$ are fixed. In particular,
the integral trace dual of $A$ is $A$ itself. The construction of the
host is a separate arithmetic question. None of the arguments below
uses the existence of column laws, comparison maps, or a later-stage
geometric point.

Write $\iota:K\hookrightarrow\mathbb C$ for the distinguished embedding,
$\theta=\sqrt{-d}$, and $\sigma$ for conjugation. The actual hypotheses
give $d\equiv7\pmod {16}$ and $\iota(\theta)=i\sqrt d$.
The trace form is linear in its first variable:
\[
 h(x,y)=\operatorname{Tr}_{E/K}(x\bar y).
                                                               \tag{PA.1}
\]
This convention fixes all transposes and conjugates below.

\paragraph{PA1: ideal lattices, area, and the discriminant.}
Let $I$ be a nonzero integral ideal of $\mathcal O_K$, and let
$\lambda\in\mathbb C^\times$. Then
$\Lambda=\lambda\iota(I)$ is a genuine rank-two complex lattice.
Indeed, $I$ is a free abelian group of rank two: it has finite index
in $\mathcal O_K$, since for $0\ne a\in I$ the full-rank multiplication
matrix of $a$ has nonzero integral determinant. For such an $a$,
$a$ and $\theta a$ belong to the rational span of $I$ and have
real-linearly independent complex images. Consequently any integral
basis of $I$ gives a real basis of $\mathbb C$. Its image of
$\mathbb Z^2$ is discrete. The additive-subgroup closure in the source
is already this additive subgroup; it is not a topological closure.

For an integral basis $(\omega_1,\omega_2)$ of $\Lambda$, the source
area, defined as an infimum of positive areas of lattice pairs, is
\[
 A_\Lambda=|\operatorname{Im}(\omega_1\bar\omega_2)|>0.
                                                               \tag{PA.2}
\]
Every other pair has area an integral determinant times (PA.2),
and the basis pair attains the least positive value. This proves
the exact area dictionary, including the absolute value.

The lattice sums
$g_2=60\sum_{0\ne\omega\in\Lambda}\omega^{-4}$ and
$g_3=140\sum_{0\ne\omega\in\Lambda}\omega^{-6}$ converge absolutely:
the number of lattice points in a disk of radius $R$ is $O((1+R)^2)$.
They therefore agree with the period-pair sums by reindexing.
Their scaling laws are $g_2(\lambda\Lambda)=\lambda^{-4}g_2(\Lambda)$
and $g_3(\lambda\Lambda)=\lambda^{-6}g_3(\Lambda)$.
Also
\[
 \Delta_\Lambda=g_2^3-27g_3^2\ne0,\qquad
 \Delta_{\lambda\Lambda}=\lambda^{-12}\Delta_\Lambda.
                                                               \tag{PA.3}
\]
Here is the standard analytic justification of nonvanishing. The
normally convergent Weierstrass series defines an even elliptic
function $\wp$ with one double pole per fundamental parallelogram.
Comparing Laurent series and cancelling poles gives
$\wp'^2=4\wp^3-g_2\wp-g_3$. The three nonzero half-period classes
are zeros of $\wp'$, by oddness and periodicity. Their $\wp$-values
are distinct: otherwise $\wp-e$ would have two distinct zeros of
order at least two, despite having just one pole of order two.
The equality of numbers of zeros and poles follows by integrating
the logarithmic derivative around a translated fundamental
parallelogram, whose opposite edges cancel. Thus the cubic has
three distinct roots, proving (PA.3). This argument also fixes the
source's normalization of $g_2,g_3,\Delta$.

\paragraph{PA2: the specified normalized period pair and half-period.}
Take $\Lambda_0=\iota(\mathcal O_K)$. Choose
$\delta^{12}=\Delta_{\Lambda_0}$, so $\delta\ne0$, and use PA1
to choose a period pair with lattice $\Lambda=\delta\Lambda_0$.
Then $\Delta_\Lambda=1$. Set $\omega_0=\delta$ and $t_0=\delta/2$.
The exact three clauses required in
\texttt{IsTwoTorsionPointUnkilled} are
\[
 2t_0\in\Lambda,\qquad
 \frac{1+i\sqrt d}{2}t_0\notin\Lambda,\qquad
 \frac{1-i\sqrt d}{2}t_0\notin\Lambda.                         \tag{PA.4}
\]
The first uses $1\in\mathcal O_K$. For either other clause,
membership and cancellation of $\delta$ would make
$(1\pm i\sqrt d)/4$ an algebraic integer. Its sum with its complex
conjugate is $1/2$, which is not an algebraic integer: a rational
algebraic integer is an integer, by the rational-root argument for
a monic integral polynomial. This proves (PA.4). In particular
$t_0\notin\Lambda$, since $(1+\theta)/2\in\mathcal O_K$; that
integrality follows from
$X^2-X+(d+1)/4$.
These are the actual period and point in
\texttt{exists\_periodPair\_lattice\_eq\_lamOf}, with no theta
nonvanishing input.

\paragraph{PA3: four copies of the trace form have a rational isometry.}
The host's embeddings indexed by $\tau$ are exactly the $K$-embeddings
of $E$. Conjugation commutes with them, since both composites become
ordinary complex conjugation after a complex embedding. Hence
$h(x,x)$ is fixed by conjugation in $K$, and is rational. If $x\ne0$,
\[
 \iota(h(x,x))=\sum_\tau |\tau(x)|^2>0.                       \tag{PA.5}
\]
Gram--Schmidt over $K$ therefore gives a basis $v_1,\ldots,v_n$
with $h(v_i,v_j)=0$ for $i\ne j$ and
$h(v_i,v_i)=a_i\in\mathbb Q_{>0}$. At each step subtract
$h(x,v_i)a_i^{-1}v_i$; positivity ensures that the norm of every
nonzero new vector is nonzero. No square root in $K$ is needed.

For each positive rational $a$, the four-squares theorem gives
$p,q,r,s\in\mathbb Q$ with $p^2+q^2+r^2+s^2=a^{-1}$.
One obtains this rational version directly by writing $a=A/B>0$,
representing $AB$ as a sum of four integer squares, and dividing
the four integers by $A$. The matrix
\[
 C=\begin{pmatrix}
 p&q&r&s\\-q&p&-s&r\\-r&s&p&-q\\-s&-r&q&p
 \end{pmatrix}
 \quad\hbox{satisfies}\quad aCC^{\mathsf T}=aC^{\mathsf T}C=1.
                                                               \tag{PA.6}
\]
Its four rows, applied to four copies of each $v_i$, give $4n$
orthonormal vectors in $E^4$. They are a basis, so their coordinate
map is a $K$-linear isometry
\[
 \gamma_0:E^4\longrightarrow K^{\{1,2,3,4\}\times\{1,\ldots,n\}},
 \qquad
 \sum_w\gamma_0(x)_w\overline{\gamma_0(y)_w}
       =\sum_{a=1}^4\operatorname{Tr}_{E/K}(x_a\bar y_a).
                                                               \tag{PA.7}
\]
This is precisely \texttt{IsoKG}; the source's equality after
embedding $K$ into $E$ is equivalent to (PA.7). The only classical
integer theorem used in (PA.6) is Lagrange's four-squares theorem,
the same \texttt{Nat.sum\_four\_squares} used in the pinned proof.

\paragraph{PA4: the complex extension is explicit, unique, and algebraic.}
Let $b_w=\gamma_0^{-1}(e_w)$ and put
$B_{(a,\tau),w}=\tau((b_w)_a)$. The trace identity gives
\[
                         B^*B=1.
                                                               \tag{PA.8}
\]
Indeed, its $(w',w)$ entry is the embedded trace pairing of
$b_w$ and $b_{w'}$. Since $B$ is square, also $BB^*=1$.
Expanding $x=\sum_w\gamma_0(x)_w b_w$ gives
\[
 (\tau(x_a))_{a,\tau}=B(\iota(\gamma_0(x)_w))_w.
                                                               \tag{PA.9}
\]
Thus $g_{\mathbb C}=B^*$ is the required complex-linear extension.
The vectors on the left of (PA.9), as $x$ varies, span the entire
complex space, because they include the columns of $B$; this also
proves uniqueness. Every entry of $B$ is algebraic, so the entries
of $g_{\mathbb C}$ are algebraic. Consequently this specific map
and its inverse preserve vectors with algebraic coordinates.
It is not necessary to assume that an arbitrary complex-linear
extension preserves such vectors.

\paragraph{PA5: separate the two primes above a split rational prime.}
Suppose $p$ is prime and either $p=2$ or $(-d/p)=1$.
There are a prime $\mathfrak p\mid p$ and $\alpha,\beta\in\mathcal O_K$
such that
\[
 \alpha\in\mathfrak p,\qquad \beta\notin\mathfrak p,
 \qquad \beta=\bar\alpha.                                    \tag{PA.10}
\]
For $p=2$ use $a=(1+\theta)/2$. Its trace is one, and its norm
$(d+1)/4$ is even. For odd $p$, choose $r\not\equiv0\pmod p$
with $r^2\equiv-d\pmod p$ and use $a=r+\theta$. Its norm is
divisible by $p$ and its trace $2r$ is prime to $p$.
In either case choose a prime ideal above $p$. Exactly one of
$a,\bar a$ belongs to it: at least one does by their product,
and both cannot by their trace. Relabel them to obtain (PA.10).
The existence of a prime above $p$ follows by taking a maximal
ideal of the nonzero finite ring $\mathcal O_K/p\mathcal O_K$;
that ring is nonzero since $\mathcal O_K$ is a free integral
module of rank two. This proof includes the dyadic case.

\paragraph{PA6: a local basis with a denominator outside one prime.}
Let $L\subset K^r$ be a full finitely generated
$\mathcal O_K$-lattice. Localize at $\mathfrak p$. The localization
is a discrete valuation ring, and $L_{\mathfrak p}$ is free of
rank $r$: successive choice of a coordinate with least valuation
and elementary elimination proves the usual basis theorem over
a discrete valuation ring. Its basis can be multiplied, vector
by vector, by elements outside $\mathfrak p$ to give vectors
$b_j\in L$ which remain a local basis. Let $g$ have these columns.
Clearing the coordinates of a finite generating family of $L$
by one $t\notin\mathfrak p$ gives
\[
       t g^{-1}L\subset\mathcal O_K^r,
       \qquad t g\mathcal O_K^r\subset L.                      \tag{PA.11}
\]
The second containment already holds with $t=1$, before taking
the common denominator. This is the exact two-sided basis
approximation consumed by the ambient construction. It does not
assert that the global lattice is a free $\mathcal O_K$-module.

\paragraph{PA7: unitary approximation at the selected split prime.}
Here is an explicit algebraic replacement for a named approximation
theorem. Use (PA.10) and put $q=\alpha/\bar\alpha$ and
$e_N=(1+q^N)^{-1}$. In the $\mathfrak p$-adic topology,
\[
       e_N\longrightarrow1,\qquad
       \bar e_N=\frac{q^N}{1+q^N}\longrightarrow0.              \tag{PA.12}
\]
In particular these expressions are defined for $N>0$.
For a matrix $A$ with $1+A$ invertible, set
$c(A)=(1+A)^{-1}(1-A)$. Direct multiplication gives
$c(c(A))=A$ and $1+c(A)=2(1+A)^{-1}$.
For any matrix $S$ set
\[
      T_N=e_NS-\bar e_NS^*,\qquad T_N^*=-T_N.
                                                               \tag{PA.13}
\]
When $1+T_N$ is invertible, $c(T_N)$ is unitary, since
$(1-T_N)^*(1-T_N)=(1+T_N)^*(1+T_N)$ and both factors commute
with their inverses. If $S=c(A)$, (PA.12) gives
$c(T_N)\to A$. Invertibility holds for all sufficiently large
$N$, by the nonzero determinant at the limit.

To approximate an arbitrary $g\in\operatorname{GL}_r(K)$,
choose a scalar $a\ne0,-1$ with $\det(a1+g)\ne0$; only finitely
many scalars are excluded. Apply the preceding construction to
$a1$ and $a^{-1}g$, and multiply the resulting unitary matrices.
Their product $u_N$ tends to $g$. Thus, for some $N$,
$g^{-1}u_N$ and $u_N^{-1}g$ have entries in
$\mathcal O_{K,\mathfrak p}$. Clear their finitely many denominators
by one $t'\notin\mathfrak p$. With $u=u_N$ we have
\[
 u^*u=1,\qquad
 t'g^{-1}u,\ t'u^{-1}g\in M_r(\mathcal O_K).                   \tag{PA.14}
\]
The argument uses a single prime topology. Complex conjugation
is not assumed continuous for that topology; (PA.12), rather
than such a false continuity assertion, is the essential step.

\paragraph{PA8: the rational denominator and the actual AtP condition.}
Suppose $L$ is self-dual for the standard Hermitian pairing.
Use PA6--PA7 and replace $L$ by $L'=u^{-1}L$. Combining
(PA.11) and (PA.14) gives, for some $t\notin\mathfrak p$,
\[
                   tL'\subset\mathcal O_K^r,
                 \qquad t\mathcal O_K^r\subset L'.            \tag{PA.15}
\]
Self-duality also gives both containments with $\bar t$.
For example, pair $x\in L'$ with $te_j$ to obtain
$\bar t x_j\in\mathcal O_K$; conversely pair $\bar t e_j$
with every $y\in L'$ and use $ty_j\in\mathcal O_K$.
Now put $a=\bar\alpha t+\alpha\bar t$. It is integral and
fixed by conjugation, hence an integer. Moreover
$a\equiv\bar\alpha t\not\equiv0\pmod{\mathfrak p}$.
Both containments in (PA.15) therefore hold with the positive
integer $s=|a|$, and $p\nmid s$. This is slightly shorter than
the generic source proof, which uses $|a\bar a|$ and only
assumes integral norms. It is valid for the actual quadratic $K$.

Apply this to
$L=\gamma_0(A^4)\subset K^{4n}$. It is finitely generated over
$\mathcal O_K$ because $\mathcal O_E$ is finite over
$\mathcal O_K$. It spans: for $0\ne b\in A$, every $x/b\in E$
has a nonzero integer multiple integral over $\mathbb Z$, so
every $x$ has an integer multiple in $A$. Finally $A^4$ is
self-dual for the sum of trace forms; testing one coordinate
at a time proves the converse implication. Isometry (PA.7)
therefore makes $L$ self-dual. For $\gamma_p=u^{-1}\gamma_0$,
\[
 s\gamma_p(A^4)\subset\mathcal O_K^{4n},\qquad
 s\gamma_p^{-1}(\mathcal O_K^{4n})\subset A^4,
 \qquad (s,p)=1.                                               \tag{PA.16}
\]
This proves both literal clauses of \texttt{AtP}, with one
denominator working for every vector. It supplies more than the
pointwise denominators required by the record.

\paragraph{PA9: all ambient record fields are now supplied.}
Let $h_K$ be the class number of $K$. For each relevant split prime
choose $\mathfrak p_p$ as in PA5, a generator $\pi_p$ of
$\mathfrak p_p^{h_K}$, and $\gamma_p$ as in PA8. The generator
exists because the ideal class has order dividing $h_K$; it is
nonzero because the ideal is nonzero. Choose
\[
 r_p^{2h_K}=\frac{\iota(\pi_p)}{\iota(\bar\pi_p)}.
                                                               \tag{PA.17}
\]
For indices not satisfying the prime/splitting condition choose
arbitrary fallback values, using the prime above two and
$\gamma_0$. This yields all eight fields of \texttt{KapIn}
including the four universally quantified specifications.
PA2, PA3, PA4, and this choice yield every field of
\texttt{VacIn}: $L$, its literal lattice equality, $\omega_0$,
the three half-period clauses, $\gamma_0$, $g_{\mathbb C}$, its
exact evaluation formula, and the correction data.

The derivative generators $c_D$ and their one common coefficient
number field are supplied by the ideal-principalization and flat
base-change argument of TL17. Use those generators together with
the \emph{same} period and vacuum data. Then
\[
                \exists J:\texttt{AmbIn}(H),\quad J.\texttt{Ok}
                                                               \tag{PA.18}
\]
holds: the two \texttt{Ok} fields are exactly
$\delta^{12}=\Delta_{\Lambda_0}$ and \texttt{IsDgen}.
No normalized theta Taylor theorem is a premise of (PA.18).
Its existence does not yet assert \texttt{StdOK}, a vacuum
four-copy identity, or any of the five \texttt{Later} stages.

\paragraph{PA10: the correction product has genuine finite support.}
The elementary determinant argument also explains the correction
factor in later vacuum formulas. Throughout this paragraph $p$
ranges over the rational primes satisfying the literal
$\texttt{SplitK}(H,p)$ predicate; the product in (PA.21) is over
that same subtype. The arbitrary fallback values of \texttt{KapIn}
outside this subtype play no role. For a fixed rational isometry
$\gamma_0$, choose a positive integer $N$ with
\[
 N\gamma_0(A^4)\subset\mathcal O_K^{4n},\qquad
 N\gamma_0^{-1}(\mathcal O_K^{4n})\subset A^4.                  \tag{PA.19}
\]
Finite generation and rational spanning give such a common $N$.
At a rational prime $p\nmid N$, both $\gamma_0$ and $\gamma_p$
identify the same localized lattice with the standard one.
Their transition and its inverse are integral at each prime of
$K$ above $p$, by (PA.16). Its determinant is therefore a unit
there. Thus the source exponent
\[
 y_p=\operatorname{ord}_{\mathfrak p_p}
                  \det(\gamma_p\gamma_0^{-1})
                                                               \tag{PA.20}
\]
vanishes unless $p\mid N$. For a fixed prime, replacing
$\gamma_p$ by another isometry satisfying \texttt{AtP} changes
the determinant by a local unit, so $y_p$ is unchanged.
This proves both \texttt{YFinite} and \texttt{YIndep}.
Each $r_p$ in (PA.17) is nonzero and algebraic, and hence
\[
                     \kappa=\prod_p r_p^{-y_p}
                     \in\overline{\mathbb Q}^{\times}.         \tag{PA.21}
\]
This is a finite product for the actual data. No inference about
support is made merely from Lean's totalized \texttt{finprod}.
The assertion concerns independence from the local isometries;
it does not identify arbitrary independent choices of the roots
$r_p$ in (PA.17).

\paragraph{PA11: a modular function with an integral unit leading term.}
We record the classical modular calculation needed for the exact
CM integrality input. Use the usual normalization
\[
 j(\tau)=\frac{E_4(\tau)^3}{\Delta_{\rm mod}(\tau)},\qquad
 q=e^{2\pi i\tau},\qquad j=q^{-1}+744+\cdots\in q^{-1}\mathbb Z[[q]].
                                                               \tag{PA.22}
\]
The source calls this function \texttt{modularJ}. It is
$1728g_2^3/(g_2^3-27g_3^2)$ for the lattice
$\mathbb Z+\mathbb Z\tau$, not $g_2^3/\Delta_\Lambda$.
For completeness, the integral expansion follows directly from
$E_4=1+240\sum_{n\ge1}\sigma_3(n)q^n$,
$E_6=1-504\sum_{n\ge1}\sigma_5(n)q^n$, and
$\Delta_{\rm mod}=(E_4^3-E_6^2)/1728=q+O(q^2)$.
The integer divisibility needed for the quotient is
$12\mid5\sigma_3(n)+7\sigma_5(n)$, obtained term by term from
$12\mid5a^3+7a^5$. Expansion then gives an integral power
series for $\Delta_{\rm mod}/q$ with constant term one; its
formal inverse is integral by recursive coefficient comparison.
The standard Eisenstein-series transformation and Fourier
formulas justify these convergent series. They follow by
reindexing their absolutely convergent lattice sums and the
partial-fraction expansion of $\pi\cot\pi z$, respectively.
Nonvanishing of $\Delta_{\rm mod}$ on the upper half-plane
follows equally from PA1 and the nonzero normalization constant.

Let $F$ be holomorphic on the upper half-plane, invariant under
$\operatorname{SL}_2(\mathbb Z)$, and suppose it has a convergent
Laurent expansion at the cusp with finitely many negative powers,
all coefficients algebraic integers, and leading coefficient an
algebraic unit $c$. Then
\[
                       F=P(j),\quad P\in\overline{\mathbb Z}[X],
                       \operatorname{lc}(P)=c.                 \tag{PA.23}
\]
To prove this, subtract $c j^m$ to cancel the leading pole
$c q^{-m}$ and continue by decreasing $m$. The remainder is
invariant, holomorphic, and bounded at the cusp, and therefore
constant: on the compactified standard fundamental domain it
extends across the cusp and across the finite elliptic quotient
charts, so the maximum principle applies. The final constant is
its integral constant Fourier coefficient. This proves (PA.23)
with all coefficient assertions. If $F(\tau_0)=0$, division by
the algebraic unit $c$ gives a monic polynomial with algebraic
integer coefficients vanishing at $j(\tau_0)$. Integrality is
transitive, so $j(\tau_0)$ is integral over $\mathbb Z$.

The same conclusion holds if the initial expansion is in
$q^{1/n}$ with an algebraic-unit leading coefficient. Raising
$F$ to the $n$th power makes its leading exponent integral;
periodicity under $\tau\mapsto\tau+1$ forces every nonintegral
Fourier exponent to have zero coefficient. All remaining
coefficients are still algebraic integers, and the leading
coefficient is $c^n$. This is the exact unit-leading-expansion
argument in the source.

\paragraph{PA12: an explicit nonsquare determinant proves CM integrality.}
Suppose $a\tau^2+b\tau+c=0$, where $a,b,c\in\mathbb Z$,
$a\ne0$, and $\operatorname{Im}\tau>0$. Put
$D=4ac-b^2>0$, $k=D+|b|$, and
\[
 M=\begin{pmatrix}k-b&-c\\a&k\end{pmatrix},\qquad
 n=\det M=k^2-bk+ac.
                                                               \tag{PA.24}
\]
Then $M\tau=\tau$ and
$4n=(2k-b)^2+D$. In particular $n>0$ and is not a square.
Indeed $m=2k-b\ge D>0$, so
$m^2<m^2+D<(m+1)^2$; the right side of $4n=m^2+D$
cannot be an integer square.

Consider the finite set of holomorphic functions $j(Az)$ as
$A$ ranges over all integral matrices of determinant $n$.
It is finite because left integral row operations reduce $A$
to a triangular matrix
$\left(\begin{smallmatrix}u&v\\0&w\end{smallmatrix}\right)$
with $u,w>0$, $uw=n$, and $0\le v<w$. Left multiplication by
$\operatorname{SL}_2(\mathbb Z)$ does not change the function.
Right multiplication permutes this finite set. Consequently
\[
             F(z)=\prod_{G\in\{j(A\,\cdot)\}}(j(z)-G(z))
                                                               \tag{PA.25}
\]
is invariant and holomorphic. Each factor has an integral
Laurent expansion in $q^{1/n}$ with an algebraic-unit leading
coefficient. In fact
$j((uz+v)/w)=\zeta_w^{-v}q^{-u/w}+O(1)$, and
$u/w\ne1$ because $n$ is nonsquare. The leading term of the
difference is therefore either $q^{-1}$ or
$-\zeta_w^{-v}q^{-u/w}$, with no leading cancellation.
All remaining coefficients lie in a cyclotomic integer ring.
The product is nonzero as a function and has a unit leading
term. The factor supplied by $M$ vanishes at $\tau$.
PA11 now proves that $j(\tau)$ is an algebraic integer.
Using the set of distinct functions, as the source does, avoids
any unjustified multiplicity count.

If a lattice $\mathbb Z+\mathbb Z\tau$ is stable under a
scalar $\alpha\notin\mathbb Z$, write
$\alpha\tau=p\tau+q$ and $\alpha=r\tau+s$ with integers
$p,q,r,s$. Then $r\ne0$ and
$r\tau^2+(s-p)\tau-q=0$. Thus the preceding argument applies.
An action by $\mathcal O_K$ supplies such an $\alpha$, since
$\mathcal O_K$ has integral rank two. This proves the literal
\texttt{isIntegral\_modularJ\_of\_ringOfIntegers} assertion.

\paragraph{PA13: the actual base $j$-value and the finite FieldOK extension.}
The base lattice in the source is $\iota(H.\mathrm{ccK})$,
where $H.\mathrm{ccK}\ne0$. By PA1 it has a period pair.
After homothety and a possible change of oriented basis it is
$\mathbb Z+\mathbb Z\tau$ with $\operatorname{Im}\tau>0$,
and is stable under $\mathcal O_K$. PA12 and the invariant
dictionary therefore give
\[
 \texttt{jLat}(\texttt{baseLat}(H))
       =1728\,g_2^3/\Delta\quad\hbox{integral over }\mathbb Z.
                                                               \tag{PA.26}
\]
This proves the exact \texttt{isIntegral\_jLat} target.

The remaining ideal in \texttt{FieldOK} is also explicit.
From $\mathfrak D_{E/K}=\mathfrak b\bar{\mathfrak b}$ and
$A=\mathfrak b^{-1}$ we have
$A\mathfrak D_{E/K}=\bar{\mathfrak b}$. Hence each
\texttt{ellLat} is the extension of a Galois conjugate of
$\bar{\mathfrak b}$. Choose $h>0$ and $\beta$ with
$\bar{\mathfrak b}^{h}=(\beta)$, and adjoin one root
$z_g^h=\varphi_0(\sigma_g\beta)$ for each of the finitely many
$g$. Also adjoin (PA.26) to the prescribed starting number field
$M$. The resulting $M'$ is finite over $\mathbb Q$ and contains
$M$. In any compatible coefficient number field containing
$M'$, the extended ideal is $(z_g)$: their $h$th powers agree,
and equality of nonzero ideals follows prime by prime from
$h v_{\mathfrak q}(I)=h v_{\mathfrak q}(J)$.
This supplies both \texttt{FieldOK} clauses with their literal
embedding compatibility and universal quantifier over later
coefficient fields. The finite ideal extension is the argument
already reconstructed in TL17; PA26 supplies its previously
retained CM-$j$ input. A fixed coefficient field may only be
used when it contains this prescribed $M'$.

\paragraph{PA14: normalized invariants and torsion coordinates are algebraic.}
Let $\Lambda=\delta\iota(I)$ be any actual curve-ideal lattice
with $\delta^{12}=\Delta_{\iota(I)}$. PA1 gives
$\Delta_\Lambda=1$, and PA12 gives integral $j(\Lambda)$.
Therefore
\[
 g_2(\Lambda)^3=j(\Lambda)/1728,\qquad
 g_3(\Lambda)^2=(g_2(\Lambda)^3-1)/27                         \tag{PA.27}
\]
show $g_2,g_3\in\overline{\mathbb Q}$. The other entries of
\texttt{DataIn} are embeddings of elements of the number field
$E$, hence algebraic as well.

If $mt\in\Lambda$, $m>0$, and $t\notin\Lambda$, then
\[
                        \wp_\Lambda(t),\ \wp'_\Lambda(t)
                              \in\overline{\mathbb Q}.         \tag{PA.28}
\]
One elementary proof keeps the curve and its normalizations
visible. The map $z\mapsto(\wp(z),\wp'(z)/2)$ identifies
$\mathbb C/\Lambda$ with the nonsingular curve
$y^2=x^3-g_2x/4-g_3/4$, carrying addition to the rational
chord-and-tangent law. The differential equation in PA1 and
the elliptic addition formula prove this identification; the
degree-two function $\wp$ is surjective, and its fibers are
the pairs $\{z,-z\}$. The $m$-torsion has exactly $m^2$ complex
points, since it is $(m^{-1}\Lambda)/\Lambda$.
Its equations and addition law are defined over
$\overline{\mathbb Q}$. If one coordinate were transcendental,
a field embedding fixing $\overline{\mathbb Q}$ could send it
to infinitely many values: send a transcendental $x$ to $x+n$,
extend to the algebraic closure of $\overline{\mathbb Q}(x)$,
and then, if desired, extend a transcendence basis to all of
$\mathbb C$. Each image would still be an $m$-torsion point,
contradicting finiteness. The second coordinate is algebraic
over the first by the displayed cubic. This proves (PA.28)
without assuming a transcendental period is algebraic.
Together, (PA.27)--(PA.28) reconstruct the actual
\texttt{admJ\_and\_torsion\_data\_algebraic} supplier.
They do not yet prove algebraicity of normalized sigma values
or Taylor coefficients at those points.

\paragraph{PA15: exact torsion-translation transfer to the even theta function.}
This last result has an explicit remaining premise. Suppose
$\Lambda$ is a genuine lattice, $s=\mathrm{sgStar}_\Lambda$
is its normalized odd theta function, and every rational-lattice
translate $U_xs$ is entire with algebraic Taylor coefficients at
zero. Here
\[
 U_xf(z)=f(z+x)\exp\!\left(-\frac{\pi z\bar x}{A_\Lambda}
                         -\frac{\pi x\bar x}{2A_\Lambda}\right).
                                                               \tag{PA.29}
\]
The ordinary analytic laws of $s$ require no arithmetic premise.
To see them, start from the normally convergent sigma product
$\sigma_\Lambda(z)=z\prod_{0\ne\omega\in\Lambda}
 (1-z/\omega)\exp(z/\omega+z^2/(2\omega^2))$.
It is odd, has derivative one at zero, and has simple zeros
exactly at the lattice. Its quasi-periods satisfy the Legendre
relation; this follows by integrating its logarithmic derivative
around a fundamental parallelogram containing its one simple
zero. For a positively oriented basis the relation is
$\eta_1\omega_2-\eta_2\omega_1=2\pi i$. It implies that there
is a complex $a$ with
$\eta_\omega=a\omega+\pi\bar\omega/A_\Lambda$ on both basis
periods, and hence on every period. Multiplying the sigma
product by $\exp(-az^2/2)$ gives the normalized $s$ and the
law $U_\omega s=\eta(\omega)s$, where the sign is positive
exactly on $2\Lambda$. These are the zero, derivative, and
translation laws used here. Their analytic normalization makes
no assertion that $a$ or a translated value is algebraic.
To identify this product with the literal chosen \texttt{sgStar},
note that any odd normalized entire function with these laws
vanishes at zero and therefore at every lattice point. Dividing
it by the product removes those possible simple poles; the
quotient is entire and periodic, hence constant on the compact
torus. The derivative-one normalization makes that constant one.

For rational-lattice points $x,y$ define
$\epsilon(x,y)=\exp(\pi(y\bar x-x\bar y)/(2A_\Lambda))$.
Expansion of exponents gives
\[
 U_xU_y=\epsilon(x,y)U_{x+y}.                                  \tag{PA.30}
\]
If $mx,ny\in\Lambda$, the signed area of this pair is an
integer multiple of $A_\Lambda$. Thus
$\epsilon(x,y)^{2mn}=1$, so these phases are nonzero algebraic
numbers. This retains the sign and the factor two of the source.

For $t\notin\Lambda$ with $2t\in\Lambda$, the source's even
theta function is $\vartheta_t=U_ts/(U_ts)(0)$. The denominator
is nonzero because the zeros of $s$ are precisely $\Lambda$;
its algebraicity follows from the stated translate premise.
Equation (PA.30) gives the identity of entire functions
\[
 U_x\vartheta_t=
       \frac{\epsilon(x,t)}{(U_ts)(0)}\,U_{x+t}s.                \tag{PA.31}
\]
Its Taylor coefficients are algebraic. For $\omega\in\Lambda$
the odd-theta law is
$U_\omega s=\eta(\omega)s$, where
$\eta(\omega)=1$ if $\omega/2\in\Lambda$ and $-1$ otherwise.
Combining it with (PA.30)--(PA.31) gives
\[
 U_{x+\omega}\vartheta_t=
 \frac{\epsilon(x+\omega,t)}{\epsilon(x,t)}
 \frac{\eta(\omega)}{\epsilon(x+t,\omega)}\,U_x\vartheta_t.
                                                               \tag{PA.32}
\]
The multiplier is a nonzero algebraic number. Entirety follows
from translation and multiplication by an exponential, and
Taylor coefficients commute with multiplication by the constant.
This is the complete stated theorem
\texttt{algebraicTaylor\_torsionTranslate\_evenTheta\_alt}; its
odd-theta translate premise remains explicit. PA14 supplies
algebraic elliptic invariants and torsion coordinates, which are
inputs to the sigma/Taylor reciprocity argument reconstructed
separately in TA; it does not replace that argument.

\paragraph{Boundary after the period construction.}
The ambient carrier, its actual local corrections, and the
CM-$j$/coefficient-field extension are now constructed from the
host. The primitive normalized odd-theta Taylor/reciprocity theorem
is supplied in the following TA section. The vacuum four-copy
identity and the exact normalized indicator and covariance
identities still need their own proofs.
The complete reference and twisted column constructions remain
conditional on the expressly stated host, index and later-stage
inputs. Nothing in PA1--PA15 supplies the independent global
spectral, metric, cohomological, or simultaneous-prime estimates.
